% ECONOMICS_PROBLEM: {"id": "C73-25", "title": "Minimizer selection by the uncorrected vanishing-common-noise game", "jel_code": "C73", "source_id": "OP-0215"}
% !TeX program = xelatex
\documentclass[11pt,a4paper]{article}
\usepackage[margin=23mm]{geometry}
\usepackage{fontspec}
\setmainfont{Latin Modern Roman}
\usepackage{amsmath,amssymb,mathtools}
\usepackage{xcolor,enumitem,booktabs,longtable,array,xurl,hyperref}
\definecolor{muted}{HTML}{53636C}
\hypersetup{colorlinks=true,urlcolor=blue,linkcolor=blue}
\setlength{\parindent}{0pt}
\setlength{\parskip}{5pt}
\newcommand{\R}{\mathbb R}
\newcommand{\E}{\mathbb E}
\newcommand{\N}{\mathbb N}
\newcommand{\ind}{\mathbf 1}
\newcommand{\Var}{\operatorname{Var}}
\newcommand{\Cov}{\operatorname{Cov}}
\newcommand{\supp}{\operatorname{supp}}
\DeclareMathOperator*{\argmin}{arg\,min}
\DeclareMathOperator*{\argmax}{arg\,max}
\newcommand{\entrymeta}[3]{{\sffamily\small\color{muted}\textbf{\texttt{#1}}\quad|\quad #2\par #3\par}}
\newcommand{\Problemref}[1]{\texttt{#1}}
\newcommand{\JELref}[2]{\texttt{#2} (JEL \texttt{#1})}
\begin{document}
\section*{Minimizer selection by the uncorrected vanishing-common-noise game}
\textbf{Problem ID:} \texttt{C73-25}\quad \textbf{JEL:} \texttt{C73}\par
% BEGIN ECONOMICS STATEMENT
\label{problem:OP-0215}
% GAME_THEORY_PROBLEM: {"local_id": "GT-085", "original_number": 85, "kind": "P", "entry_kind": "statement_only", "published": false, "review_required": true, "category": "game_theory"}
\label{game:GT-085}
{\sffamily\small\color{muted}
{\small\textbf{GT-085} | Mean-field games | P: published question with an essential regularization boundary\par}
\par}
\subsection*{Definitions and mathematical statement}
Fix $E=\{1,\ldots,d\}$, $d\geq3$, $T>0$, $m_0\in\operatorname{int}\Delta_d$, where $\Delta_d=\{m\geq0:\sum_i m_i=1\}$. Let $\mathcal F,\mathcal G$ be $C^\infty$ potentials on a neighborhood of the simplex, with $f_i=\partial_i\mathcal F$, $g_i=\partial_i\mathcal G$. Write $r_+=\max(r,0)$ and
\[
 H_i(U)=\tfrac12\sum_{j\ne i}(U_i-U_j)_+^2.
\]
For $\varepsilon>0$, an admissible boundary repulsion is a smooth nonnegative function $\phi_\varepsilon$ with $\phi_\varepsilon\to0$ locally uniformly on $(0,1]$, for which the following penalized master equation has a unique classical solution $U^\varepsilon$. Require $U^\varepsilon$ to be continuous on $[0,T]\times\Delta_d$, $C^1$ in time and $C^2$ in the simplex interior, with bounded value and first derivatives and second derivatives bounded on interior compact sets; require the population SDE below to have a unique interior-valued solution. These are explicit admissibility conditions, ensured in the source by sufficiently strong repulsion at zero. Define
\[
 b_i^\varepsilon(m,U)=\sum_{j\ne i}
 \left[m_j\{\phi_\varepsilon(m_i)+(U_j-U_i)_+\}
       -m_i\{\phi_\varepsilon(m_j)+(U_i-U_j)_+\}\right].
\]
The master equation, with the original costs $f,g$ kept unchanged, is
\[
 \begin{split}
 0={}&\partial_tU_i^\varepsilon
 +\varepsilon\sum_{j,k}(m_j\delta_{jk}-m_jm_k)\partial_{m_jm_k}^2U_i^\varepsilon
 +\sum_j b_j^\varepsilon(m,U^\varepsilon)\partial_{m_j}U_i^\varepsilon\\
 &+2\varepsilon\sum_jm_j(\partial_{m_i}U_i^\varepsilon-\partial_{m_j}U_i^\varepsilon)
 -H_i(U^\varepsilon)+f_i(m)
 +\sum_{j\ne i}\phi_\varepsilon(m_j)(U_j^\varepsilon-U_i^\varepsilon),\\
 U_i^\varepsilon(T,m)={}&g_i(m).
 \end{split}
\]
$\delta_{jk}$ is the Kronecker delta. All derivatives are intrinsic tangent derivatives, equivalently derivatives of an extension; the tangent contractions displayed here are independent of the extension. In particular, this is a Kimura diffusion on the simplex, not an independent additive Brownian noise for each coordinate.

With independent Brownian motions $(W^{ij})_{i,j\in E}$, the selected population law solves
\[
 dm_i^\varepsilon=b_i^\varepsilon(m^\varepsilon,U^\varepsilon(t,m^\varepsilon))dt
 +\sqrt\varepsilon\sum_j\sqrt{m_i^\varepsilon m_j^\varepsilon}(dW^{ij}-dW^{ji}),
 \qquad m^\varepsilon(0)=m_0.
\]
Define $\mathcal M_{T,m_0}$ as the law paths in global minimizers of
\[
 \int_0^T\left[\tfrac12\sum_i\sum_{j\ne i}m_i a_{ij}^2+\mathcal F(m)\right]dt
 +\mathcal G(m(T)),\qquad
 \dot m_i=\sum_{j\ne i}(m_ja_{ji}-m_ia_{ij}),\quad a_{ij}\geq0,
\]
with $m(0)=m_0$, absolutely continuous $m\in\Delta_d$, and measurable finite-energy rates. For admissible repulsion families and convergent subsequences
\[
 \mathcal L(m^{\varepsilon_n})\Longrightarrow Q
 \quad\text{in }C([0,T],\Delta_d),\qquad \varepsilon_n\downarrow0,
\]
determine whether $Q(\mathcal M_{T,m_0})=1$. A full statement of any proposed theorem must specify its boundary-layer scale and the estimates ensuring tightness and removal of the repulsion in the limit.
\subsection*{Short English statement}
Does the exact noisy game, with only the boundary repulsion and no extra cost correction, select deterministic potential minimizers as its common noise vanishes?
\subsection*{Variants and scope boundaries}
The restriction to unchanged $f$ is essential. The source's Theorem~4.5 obtains minimizing selection after introducing additional costs $f^\varepsilon$ that restore a stochastic potential structure; that is a known positive result. Its approximate-equilibrium conclusion does not settle the exact uncorrected question. The finite-state Wright--Fisher formulation also uses a weighted representative-player interpretation, not an ordinary unweighted conditional jump law.
\subsection*{Source relationship and status}
The uncorrected question appears explicitly before Theorem~4.5. The theorem's cost-corrected variant is already resolved in the cited source; its different quantifiers must be retained.
\subsection*{References}
\begingroup\footnotesize\raggedright
{\footnotesize\raggedright
\textbf{[S1]} Cardaliaguet and Delarue. \emph{Selected topics in mean field games}, ICM2022, vol.~5; DOI: 10.4171/ICM2022/17.\par
\textit{Locators:} (4.14), p.~3686, penalized master equation; (4.15), Theorem~4.3, p.~3687; Definition~4.1, pp.~3687--3688; Section~4.3, ``Selection of equilibria,'' p.~3690, exact uncorrected question; Theorem~4.5, pp.~3691--3692, cost-corrected positive result.\par
\url{https://ems.press/content/book-chapter-files/33265}\par\smallskip
\textbf{[S2]} Alekos Cecchin and Fran\c{c}ois Delarue. \emph{Selection by vanishing common noise for potential finite state mean field games}. Communications in Partial Differential Equations 47(1) (2022), 89--168; DOI: 10.1080/03605302.2021.1955256.\par
\textit{Source role:} primary paper for the corrected selection and intrinsic weak uniqueness results cited as [43] by the chapter; publication date differs from the chapter's early citation.\par
\url{https://arxiv.org/abs/2005.12153}\par}
\par\endgroup

\subsection*{Common conventions: Source mathematical conventions}

\textbf{Finite games.} For a finite set $A$, $\Delta(A)$ is its probability
simplex. Players choose independent mixed strategies in Nash equilibrium;
correlation is allowed only when explicitly part of the solution concept.
In a game with payoff functions $u_i$ and strategy profile $x$, an additive
$\varepsilon$ Nash equilibrium satisfies
\[
 u_i(x)\geq u_i(x_i',x_{-i})-\varepsilon
 \quad\text{for every player $i$ and every admissible deviation $x_i'$}.
\]
For numerical approximation, payoffs are normalized as locally specified.
An $\varepsilon$ payoff residual is distinct from distance at most
$\varepsilon$ to the set of exact equilibria. Point convergence, convergence
of empirical play, and convergence in distance to an equilibrium set are
also distinct.

\textbf{Computational representations.} Unless stated otherwise, a complexity
question takes rational data with binary encoding; polynomial time is measured
in total bit length. A fixed-accuracy polynomial algorithm is not necessarily
polynomial in $1/\varepsilon$, and neither is necessarily polynomial in
$\log(1/\varepsilon)$. Succinct games and value, demand, or sampling oracles
must declare their interfaces. Exact real solutions require an explicit
representation; an unspecified list of arbitrary real numbers is not a
finite computational output. A claimed algorithmic barrier is meaningful
only for its specified model and reductions.

\textbf{Dynamic games.} Histories, information, observation, and allowable
randomization are part of the model. A stationary strategy depends only on
the current state; a positional strategy in a deterministic graph game
chooses one action at each controlled vertex. Nash equilibrium at an initial
state and equilibrium after every history are different requirements.
An expectation of a pathwise $\liminf$ payoff, a $\liminf$ of expected averages,
a limiting expected average, and a uniform finite-horizon equilibrium are
different criteria. None is silently substituted for another.

\textbf{Allocation and mechanisms.} A complete allocation assigns every item
exactly once unless otherwise stated. Zero-valued goods, negative values,
capacity constraints, feasibility, lotteries, and copies of goods affect
fairness definitions and are specified locally. Dominant-strategy incentive
compatibility, Bayesian incentive compatibility, universal truthfulness,
and truthfulness in expectation impose different inequalities. Whenever
an objective ratio has zero denominator, its convention is stated or those
instances are excluded.

\textbf{Infinite-dimensional models.} Expectations, integrals, stochastic
equations, and optimization objectives require their local regularity,
integrability, filtrations, and admissible domains. A backward--forward
mean-field system is a boundary-value problem; it is not an initial-value
semiflow without an additional construction. Long-time, large-population,
and vanishing-noise limits require an order of limits and a convergence
topology. If a minimum need not be attained, the objective is its infimum
or supremum and the approximation requirement is stated separately.
% END ECONOMICS STATEMENT
\end{document}
